\begin{recipe}{Asperges à la flamande}
\kicker[Hoofdgerecht,Vlees,Vlaams]{Hoofdgerecht · vlees · Vlaams}
\index[register]{Asperges}
\index[register]{Krieltjes}
\index[register]{Ham}

\meta{30 minuten \dvd Voor 4 personen}

\lettrine{A}{sperges} zijn hier op de Brabantse Wal een traditie. In het
seizoen koop je ze het liefst vers bij de boer, bijvoorbeeld bij
boerderijcamping De Meet, waar ze meteen voor je geschild worden.

\blockrule{Ingrediënten}
\begin{ingredients}
  \ing{1 kg}{verse witte asperges}\index[register]{Asperges}
  \ing{1 kg}{krieltjes}\index[register]{Krieltjes}
  \ing{6}{middelgrote scharreleieren}
  \ing{10 g}{verse krulpeterselie}
  \ing{125 g}{ongezouten roomboter}
  \ing{8}{plakjes achterham}\index[register]{Ham}
\end{ingredients}

\blockrule{Bereiding}
\tip{Bij ons hoort de traditie van de drie A's erbij: aardappelen, asperges
en aardbeien. Sluit af met aardbeien en een likje vanille yoghurt als
toetje, als het aspergeseizoen dat toelaat.}
\begin{steps}
  \item Snijd de peterselie alvast fijn, zodat je daar straks niet meer
        naar om hoeft te kijken.
  \item Schil de asperges met een dunschiller rondom, vanaf vlak onder het
        kopje naar beneden, en verwijder het houtachtige onderste stuk.
        Halveer de grotere krieltjes.
  \item Zet drie pannen met ruim water op, eventueel gezouten: een
        aspergepan voor de asperges, een gewone pan voor de krieltjes en
        een steelpannetje voor de eieren. Breng alle drie aan de kook.
  \item Kook de asperges in de aspergepan 10 minuten beetgaar, de
        krieltjes in de andere pan 8 tot 10 minuten gaar, en de eieren in
        het steelpannetje 8 minuten hard. Zet ze tegelijk op, dan zijn ze
        ook tegelijk klaar.
  \item Smelt ondertussen de boter in een kleine (steel)pan op laag vuur,
        vlak voor je gaat opdienen, zodat hij niet alvast weer afkoelt.
  \item Giet de eieren af zodra ze gaar zijn, koel ze onder koud water en
        pel ze. Prak ze grof met een vork.
  \item Giet de asperges en de krieltjes af en verdeel ze over de borden.
        Verdeel achtereenvolgens het geprakte ei, de gesmolten boter en de
        peterselie over de asperges. Verdeel de ham in mooie roosjes en
        leg ze ernaast.
\end{steps}
\end{recipe}
